\begin{defi}\label{def.classAndInstance}
 The type of an object is itself an object, and it is called a \emph{class}.
 An object whose type is the class $C$ is an \emph{instance} of $C$:
 \codehl{type(x)} returns the class of which \codehl{x} is an instance,
 and calling a class creates a new instance of it.
\end{defi}

\codehl{dict} is a class, and a side of the book is an instance of it.
The statement \codehl{class} defines a new one:
the names bound in its body, most of them to functions,
become the \emph{attributes} of the class, written \codehl{C.name}.
An instance has attributes of its own, written \codehl{x.name}.
They are created by assignment,
and \codehl{vars(x)} returns them as a dictionary.

The accumulator of the fold \eqref{eq.bookFold} is the class \codehl{AggregateBook}.
When \codehl{AggregateBook(strict)} is evaluated,
Python creates an instance with no attributes
and calls the function \codehl{\_\_init\_\_} of the class,
with the instance as first argument and \codehl{strict} after it.
That function is the \emph{constructor}, in Listing \ref{listing.aggregateInit}:
by assigning to the attributes of its first parameter it gives the instance its state,
here two empty sides, the flag \codehl{strict} of Section \ref{sec.foldImplementation},
and three counts of what the last message traded.
The first parameter is named \codehl{self} by convention,
and only its position makes it the instance.
The annotations of the listing, \codehl{strict: bool} and those after it,
state the types that are intended;
Python records them and does not check them.

\begin{lstlisting}[caption={\protect\codehl{unito26.lob.orderbook.AggregateBook.\_\_init\_\_}.}, label=listing.aggregateInit]
    def __init__(self, strict: bool = False):
        self.bids: dict[int, int] = {}
        self.asks: dict[int, int] = {}
        self.strict = strict
        self.last_fill_count = 0
        self.last_traded_size = 0
        self.last_traded_value = 0
\end{lstlisting}

Two calls of the class create two books,
each with its own six attributes,
and a write to one leaves the other as it was.
A name refers to a book, and not to the configuration the book is in:
in the fold \eqref{eq.bookFold} the configurations $\mathcal{B}_0, \mathcal{B}_1, \dots$
are the successive values of one object.
So a list to which the book is appended after every message holds one object many times,
and every entry shows the last configuration.
Whatever records the sequence stores copies,
and \codehl{copy} returns a new book in the configuration of the old one.
Listing \ref{listing.twoBooks} checks these statements.

\begin{lstlisting}[caption={Two instances of one class, and a copy.}, label=listing.twoBooks]
from unito26.lob.messages import BUY
from unito26.lob.orderbook import AggregateBook

book, other = AggregateBook(), AggregateBook()
assert type(book) is type(other) is AggregateBook
assert book is not other
assert set(vars(book)) == {
    "bids", "asks", "strict",
    "last_fill_count", "last_traded_size", "last_traded_value",
}

book.set_size(BUY, 99, 100)
assert book.bids == {99: 100} and other.bids == {}

states = [book, book.copy()]
book.set_size(BUY, 99, 0)
assert states[0] is book and states[0].bids == {}
assert states[1].bids == {99: 100}
\end{lstlisting}

\begin{remark}\label{remark.misspeltAttribute}
 The assignment \codehl{x.name = value} creates the attribute if \codehl{x} does not have it.
 On a book, \codehl{book.bid = \{\}} raises no error:
 it gives the book a third dictionary, which no function of the class reads.
\end{remark}
